\documentclass[11pt,a4paper]{article}
\usepackage[utf8]{inputenc}
\usepackage[spanish]{babel}
\usepackage{amsmath}
\usepackage{amsfonts}
\usepackage{amssymb}
\usepackage{graphicx}
\usepackage{geometry}
\geometry{left=2.5cm,right=2.5cm,top=2.5cm,bottom=2.5cm}

\title{\textbf{ANÁLISIS DE CASO 6: ZARA (INDITEX) - MODA RÁPIDA}\\ Curso: Gestión de Operaciones (ICS3213)}
\author{Integrantes: César Meneses, Fabián Ortega}
\date{29 de Junio de 2026}

\begin{document}

\maketitle

\section*{Pregunta 1}

\subsection*{a) Estrategia Competitiva y Ventaja Sustentable}
Zara sigue una estrategia genérica de diferenciación basada en la agilidad y respuesta rápida (\textit{Quick Response}), complementada con precios accesibles. Lo que fundamenta esto es su propuesta de valor orientada a ofrecer las últimas tendencias de la moda al mercado masivo en un tiempo récord de 4 a 5 semanas, creando un efecto de escasez artificial a través de una alta rotación de inventarios. 
Su ventaja competitiva radica en la integración vertical hacia atrás y su extrema flexibilidad en la cadena de suministro. A diferencia de competidores que buscan el menor costo fabril en Asia enfrentando tiempos de ciclo de 6 meses, Zara asume mayores costos de producción en Europa para minimizar este tiempo de ciclo. 
Esta ventaja es sustentable en el tiempo debido a las altas barreras de entrada operacionales. Replicar la infraestructura de tecnología, el centro de distribución dinámico de Arteixo y la cultura de manufactura celular Just-In-Time requiere décadas de inversión en capital y aprendizaje organizacional que los competidores tradicionales no pueden imitar fácilmente a corto plazo.

\subsection*{b) Segmento de Clientes, Propuesta de Valor y Competencia}
Aunque abarca hombres, mujeres y niños, su segmento principal son mujeres de 18 a 34 años de edad, de ingresos medios a altos, sumamente conscientes de las tendencias de la moda y que habitan en zonas urbanas densas. El cliente valora la novedad continua de las colecciones, el diseño a precios razonables y el entorno atractivo de la tienda, el cual induce una percepción de exclusividad y urgencia.
Sus principales competidores globales son H\&M, enfocada en bajos precios y moda rápida mediante externalización eficiente; The Gap, enfocada en básicos de calidad y sufriendo por largos ciclos; y Benetton, estructurada en red de franquicias y prendas de punto con alta inversión en teñido rápido, pero con ciclos más largos de pedido.

\subsection*{c) Cadena de Abastecimiento y Rol de Proveedores}
La cadena de abastecimiento de Zara es centralizada, híbrida y vertical. Todo fluye a través del mega-centro en Arteixo. Las fábricas propias manufacturan el 40\% del volumen, enfocándose en los artículos más inciertos para aportar flexibilidad instantánea. Los proveedores de proximidad en España, Portugal y el norte de África abastecen costura externa, permitiendo reabastecimiento en 2 semanas. Los proveedores de Asia proveen artículos básicos de baja incertidumbre para reducir costos mediante economías de escala. Adicionalmente, la subsidiaria Comditel compra el 50\% de la tela sin teñir, aplicando la estrategia de postergación para retrasar la decisión del color hasta tener datos certeros de demanda.
El flujo de información constante mediante dispositivos en tiendas elimina el efecto látigo, mientras que los lotes pequeños y los despachos bi-semanales minimizan la necesidad de bodegaje local y liquidaciones, situando los descuentos entre 15\% y 20\% versus el 30\% a 40\% de la competencia.

\newpage

\section*{Pregunta 2}

\subsection*{a) Crecimiento Internacional: Restricciones de la Cadena de Suministro}
La cadena actual permite el crecimiento dado que el modelo de operaciones está altamente estandarizado. La capacidad de consolidar envíos y realizar rutas aéreas permite llegar a Medio Oriente o América rápidamente manteniendo el efecto escasez. 
Sin embargo, el esquema actual genera antieconomías de escala. La hipercentralización en Arteixo tiene un límite volumétrico, y el sobrecosto de flete aéreo a América o Asia merma la rentabilidad respecto a las ventas europeas. Los problemas prácticos incluyen husos horarios que complican la logística en tiempo real, saturación de envíos y barreras aduaneras.
Para solucionar esto, yo modificaría la cadena descentralizando nodos de distribución por macro-regiones geográficas para artículos básicos y utilizaría el diseño de \textit{mancha de petróleo} local para prorratear costos logísticos. Inditex de hecho tomó la decisión de replicar Arteixo en Zaragoza para aliviar este cuello de botella central.

\subsection*{b) Métricas de Proveedores y Logística (KPIs)}
Para los proveedores se utilizarían tres indicadores clave. Primero, el tiempo de ciclo de cumplimiento, que mide la agilidad de los contratistas y fundamenta la estrategia de respuesta rápida. Segundo, la tasa de entregas a tiempo y completas, ya que en una red Just-in-Time un quiebre en la entrega genera ventas perdidas irremediables. Tercero, el porcentaje de defectos de calidad, para asegurar que la velocidad fabril no merme la calidad exigida por el segmento de clientes.
Para la logística se utilizarían otros tres indicadores. Primero, los días de permanencia en el centro de distribución, garantizando que el nodo actúe como tránsito y no como bodega muerta. Segundo, el costo de transporte por prenda como porcentaje de las ventas, fundamental para monitorear las ineficiencias al exportar vía aérea intercontinental. Tercero, la rotación de inventario en tienda, que refleja directamente el impacto del lote pequeño bi-semanal y asegura el clima de escasez en el punto de venta.

\subsection*{c) Internacionalización de Empresas Chilenas: El caso de Falabella}
Falabella S.A. es el ecosistema de \textit{retail} más grande de Chile. Su internacionalización comenzó en los años noventa hacia Argentina, seguido de Perú, Colombia y adquisiciones regionales como Linio.
El proceso en Perú y Colombia fue un rotundo éxito, sustentado por una alta afinidad cultural y la exportación de un modelo de negocio cruzado entre venta de bienes y su negocio financiero CMR, logrando integrarse con proveedores locales y construir centros de distribución eficaces. 
Por el contrario, el mercado en Argentina terminó en fracaso, llevando al cierre de sus operaciones. Las causas radican en las severas restricciones del gobierno a las importaciones, el control cambiario y la inflación, que erosionaron los márgenes y dificultaron el abastecimiento.
Al internacionalizar, las operaciones deben modelarse no sólo por la demanda potencial, sino por el riesgo macroeconómico de la cadena de suministro. Las barreras arancelarias, la dificultad de remitir inventarios internacionales de forma expedita y las diferencias socioculturales determinan si una ventaja competitiva local es extrapolable a nivel multinacional.

\end{document}
